\section{Conclusion}
\label{sec:conclusion}

Reading the data before choosing a model paid off more than any single technique.
Tiny objects called for a dense stride-8 detector rather than an adapted stride-16 encoder; repeated scenes called for a
holdout that mirrors val and for augmentations that keep scene geometry; the shared class of an image and its boxes
allowed a free cross-task decoding.
The result is a 26\,M-parameter single model with $\Delta_m$ +45.9 on val and +48.7 on test, over ten times the best
in-budget baseline, trained in under 10 GPU-hours.
We also learned our protocol's limits: the holdout predicts the sign of a change better than its size, seed noise
rivals most effects, and bridges and playgrounds, tiny objects in cluttered scenes with few repeats, remain unsolved.
They, and the 11-point mAP gap to the best teams, are where we would go next: condition the detection head on the
predicted scene class, so that it expects bridges in bridge scenes; give the stride-8 features more context, for
example with attention; and mine look-alike structures that carry no box as hard negatives.
